\subsection{Filtered-data ablation, kept as a negative protocol lesson}
\label{sec:filtered-ablation}
Before the full-record run, study 13 used a 150-residue cap matching the
UniProt model and dropped nine long sequences from the Feb2020 corpus.
It also removed the unknown residue X from two decoy strings before
encoding. On its 4{,}033 changed records, its stratified 10-fold CV mean
was ACC $0.9055$, MCC $0.8116$, and AUC $0.9634$; these are measured
values, but they are \emph{not} a fair result on the published 4{,}042-entry
benchmark. The full-record study 14 retains all sequences and represents
unknown positions as zero feature rows. This distinction matters: a
benchmark "improvement" can be manufactured accidentally by dropping
harder long cases. We preserve the ablation and its source code instead
of overwriting it after finding the mismatch. Both our random-split runs
still contain exact 10-mer cross-fold overlap, so neither supplants the
homology-aware result in Section~\ref{sec:results}.
